\documentclass[11pt]{article}
\usepackage[margin=1in]{geometry}
\usepackage{amsmath,amssymb,booktabs,longtable,setspace,caption,hyperref,array}
\usepackage[T1]{fontenc}
\usepackage[utf8]{inputenc}
\usepackage{mathpazo}
\usepackage{xurl}
\setlength{\parindent}{0pt}
\setlength{\parskip}{0.55em}
\captionsetup{font=small,skip=6pt}
\hypersetup{
  colorlinks=true,
  linkcolor=black,
  citecolor=black,
  urlcolor=blue
}

\title{Construction of a Quarterly Panel of Real GDP per Worker\\
in Constant 2015 USD at Market Exchange Rates}
\author{}
\date{10 September 2026}

\begin{document}
\maketitle

\begin{abstract}
\noindent
This note documents the construction of an unbalanced country--quarter panel of
\textbf{real GDP per worker} (employed persons, not population). The series is
expressed in a common international unit: constant 2015 US dollars at 2015
period-average market exchange rates. It is not PPP. Exchange rates enter
\textbf{once}, in the base year, so the time path of each country is the path of
real GDP in local currency and is not contaminated by quarterly bilateral-dollar
volatility. Employment is official ILOSTAT quarterly data where available, and
linearly interpolated official annual data otherwise. The current vintage has
101 countries and 7{,}450 quarterly observations, 1950Q1--2026Q2. The panel,
raw extracts, and code live entirely in this folder; no series from
elsewhere in the repository are used.
\end{abstract}

\tableofcontents
\clearpage

\section{Purpose and target series}

The object of interest is labour productivity at quarterly frequency, for as
many economies as national accounts and labour-force data will support:

\begin{quote}
annualized real GDP per \emph{employed person}, in constant 2015 USD at 2015
market exchange rates.
\end{quote}

Three design constraints follow from that sentence.

\begin{enumerate}
\item \textbf{Per worker, not per capita.} The denominator is employment
(persons), not mid-year population. Using population would mix participation,
demography, and productivity.
\item \textbf{Common units, not PPP.} Levels must be comparable across
countries in the same dollar vintage. Purchasing-power parities are not used.
The converter is the 2015 period-average market exchange rate (domestic
currency per USD).
\item \textbf{No quarterly FX in the volume path.} Converting each quarter with
that quarter's exchange rate would put dollar--local-currency swings into
measured productivity. Real GDP is therefore taken in local currency at
constant prices, turned into a 2015 volume index, scaled by 2015 current-price
GDP, and divided by the \emph{2015} average exchange rate only.
\end{enumerate}

The delivered file is
\texttt{output/gdp\_per\_worker\_q.csv}.
Required columns are \texttt{country} (ISO 3166-1 alpha-3), \texttt{period}
(\texttt{YYYY-Qn}), \texttt{gdp\_per\_worker}, and \texttt{metadata}. Additional
columns are audit extras (country name, year, quarter, calendar time,
annualized real GDP in 2015 USD, employment in persons, SA/NSA flag, employment
source).

\section{Construction formula}

For country $i$ and quarter $t$,
\begin{equation}
\label{eq:usd}
Y^{\mathrm{USD}}_{i,t}
=
\left(\frac{Y^{r,q}_{i,t}\times 4}{Y^{r,A}_{i,2015}}\right)
\times
\left(\frac{Y^{n,A}_{i,2015}}{E_{i,2015}}\right),
\qquad
\text{gdp per worker}_{i,t}
=
\frac{Y^{\mathrm{USD}}_{i,t}}{L_{i,t}}.
\end{equation}

\begin{table}[h]
\centering
\caption{Symbols in~\eqref{eq:usd}.}
\begin{tabular}{>{\raggedright\arraybackslash}p{0.18\textwidth}
                >{\raggedright\arraybackslash}p{0.46\textwidth}
                >{\raggedright\arraybackslash}p{0.28\textwidth}}
\toprule
Symbol & Meaning & Source \\
\midrule
$Y^{r,q}_{i,t}$ &
Real GDP, domestic currency, \textbf{quarterly flow} (not SAAR) &
IMF QNEA \texttt{B1GQ}, constant prices, \texttt{XDC} \\
$Y^{r,A}_{i,2015}$ &
Real GDP, domestic currency, 2015 annual &
IMF ANEA \texttt{B1GQ} constant \texttt{XDC};
fallback $=$ sum of four 2015 QNEA quarters \\
$Y^{n,A}_{i,2015}$ &
Current-price GDP, domestic currency, 2015 annual &
IMF ANEA \texttt{B1GQ} current \texttt{XDC} \\
$E_{i,2015}$ &
2015 period-average domestic currency per USD &
IMF ER \texttt{XDC\_USD} / \texttt{PA\_RT} / annual \\
$L_{i,t}$ &
Employed persons &
ILOSTAT (Section~\ref{sec:emp}) \\
$\times 4$ &
Quarterly flow $\to$ annualized flow &
Unit conversion, not seasonal adjustment \\
\bottomrule
\end{tabular}
\end{table}

\paragraph{Why 2015.}
Chain-linked national-accounts volumes have a country-specific reference year
(the US QNEA extract is referenced to 2017). Dividing each quarter by that
country's 2015 real annual total puts every country on a $2015=1$ volume index.
Multiplying by 2015 current-price GDP converts the index back into 2015-price
local currency. Dividing by the 2015 average FX then puts every country in the
same international unit.

\paragraph{Why multiply by 4.}
IMF QNEA constant-price GDP is a quarterly \emph{flow}, not a seasonally
adjusted annual rate. For the United States, 2015Q1 is about
$4.67\times 10^{12}$; the four 2015 quarters \emph{sum} to the ANEA 2015 annual
total ($\approx 18.80\times 10^{12}$). Annualizing by 4 makes GDP/employment
``output per worker per year.''

\paragraph{Identities (tested).}
If the four 2015 real quarters sum to the 2015 annual real total, then 2015
annualized $Y^{\mathrm{USD}}$ equals 2015 current-price GDP divided by
$E_{i,2015}$. Doubling real LCU doubles $Y^{\mathrm{USD}}$. Doubling the
base-year FX halves $Y^{\mathrm{USD}}$. The shipped rebase function has no
quarterly-FX argument.

\paragraph{United States, 2015.}
ANEA 2015 current-price GDP is $\$18.295$ trillion; $E_{\mathrm{US},2015}=1$.
ILO official quarterly employment is about 147--150 million persons. The
constructed series is therefore about $\$122{,}500$--$123{,}500$ per worker in
each 2015 quarter, which is the annualized identity above.

\section{Sources and exact extracts}
\label{sec:sources}

Extracts are cached under \texttt{raw/} by \texttt{download\_sources.py}.
Re-run with \texttt{--force} to refresh. All IMF queries use SDMX~3.0,
\texttt{Accept: text/csv}. The wildcard \texttt{*} is all countries.

\subsection{IMF National Economic Accounts}

\paragraph{Quarterly (QNEA).}
Dataset: \url{https://data.imf.org/en/datasets/IMF.STA:QNEA}.
Key order:
\path{COUNTRY.INDICATOR.PRICE_TYPE.S_ADJUSTMENT.TYPE_OF_TRANSFORMATION.FREQUENCY}.
Indicator \texttt{B1GQ} is GDP. \texttt{Q} in the price slot is constant
prices; \texttt{XDC} is domestic currency; the last \texttt{Q} is quarterly
frequency.

\begin{itemize}
\item SA volumes:
\url{https://api.imf.org/external/sdmx/3.0/data/dataflow/IMF.STA/QNEA/~/*.B1GQ.Q.SA.XDC.Q}
\item NSA volumes:
\url{https://api.imf.org/external/sdmx/3.0/data/dataflow/IMF.STA/QNEA/~/*.B1GQ.Q.NSA.XDC.Q}
\end{itemize}

\paragraph{Annual (ANEA).}
Dataset: \url{https://data.imf.org/en/datasets/IMF.STA:ANEA}.
Key order:
\path{COUNTRY.INDICATOR.PRICE_TYPE.TYPE_OF_TRANSFORMATION.FREQUENCY}.

\begin{itemize}
\item Constant-price GDP, LCU, annual:
\url{https://api.imf.org/external/sdmx/3.0/data/dataflow/IMF.STA/ANEA/~/*.B1GQ.Q.XDC.A}
\item Current-price GDP, LCU, annual:
\url{https://api.imf.org/external/sdmx/3.0/data/dataflow/IMF.STA/ANEA/~/*.B1GQ.V.XDC.A}
\end{itemize}

IMF \texttt{SCALE} is a power-of-ten unit multiplier
($v\times 10^{\mathrm{SCALE}}$). The GDP extracts used here have
\texttt{SCALE}=0 (units of 1 domestic-currency unit). Missing scale is treated
as 0. Staff projections and imputations
(\texttt{DERIVATION\_TYPE} in $\{\texttt{SP},\texttt{SE},\texttt{SEME},\texttt{SEHI}\}$)
are dropped.

ANEA \texttt{TIME\_PERIOD} arrives as float (\texttt{2015.0}) when pandas reads
the CSV. The parser accepts \texttt{2015}, \texttt{2015.0}, and integer 2015,
and rejects \texttt{nan}.

\subsection{IMF Exchange Rates}

Dataset: \url{https://data.imf.org/en/datasets/IMF.STA:ER}.
Key order:
\path{COUNTRY.INDICATOR.TYPE_OF_TRANSFORMATION.FREQUENCY}.
\texttt{XDC\_USD} is domestic currency per US dollar; \texttt{PA\_RT} is period
average; frequency is annual.

\url{https://api.imf.org/external/sdmx/3.0/data/dataflow/IMF.STA/ER/~/*.XDC_USD.PA_RT.A}

Country names come from IMF codelist \texttt{CL\_COUNTRY}:
\url{https://api.imf.org/external/sdmx/3.0/structure/codelist/IMF/CL_COUNTRY}.

\subsection{ILOSTAT employment}

Bulk CSVs via the ILO Rplumber indicator service
(\url{https://ilostat.ilo.org/data/}):

\begin{itemize}
\item Official, quarterly:
\url{https://rplumber.ilo.org/data/indicator/?id=EMP_TEMP_SEX_AGE_NB_Q}
\item Official, annual:
\url{https://rplumber.ilo.org/data/indicator/?id=EMP_TEMP_SEX_AGE_NB_A}
\item ILO modelled estimates, annual (downloaded as a last-resort source; not
used in the current vintage, because every panel country had official data):
\url{https://rplumber.ilo.org/data/indicator/?id=EMP_2EMP_SEX_AGE_NB_A}
\end{itemize}

ILO time labels are \texttt{YYYYQn} (no hyphen). IMF labels are
\texttt{YYYY-Qn}. Both parse to the same year and quarter.

\section{Seasonal adjustment of GDP}

QNEA publishes both \texttt{SA} and \texttt{NSA}. IMF SA and NSA GDP are
never spliced inside a country:

\begin{itemize}
\item If a country has any SA constant-price GDP, all of that country's GDP
observations come from SA.
\item Otherwise GDP is NSA. Optionally, productivity is then seasonally
adjusted with statsmodels SARIMAX when the spell is long enough
(\texttt{sarimax\_nsa=True} / \texttt{--sarimax-nsa}; \textbf{off by default}).
\end{itemize}

\paragraph{SARIMAX on NSA countries (optional, default off).}
NSA GDP, and typically NSA employment, leave seasonality in the productivity
ratio. When the flag is on, and there are at least 12 observations covering
all four calendar quarters, we fit
\begin{equation}
\log(\text{gdp per worker})_{t}
=
a + g t + \gamma_{q(t)} + u_{t},
\qquad
u_{t}\sim\mathrm{ARMA}(p,q),
\end{equation}
as \texttt{SARIMAX(log $y$, order=$(p,0,q)$, trend=ct, exog=Q2--Q4 dummies)}.
$(p,q)$ is chosen by AIC. \texttt{d=0} so dummy coefficients are
\emph{level} seasonal factors (differencing would turn dummies into pulses).
The four $\gamma$s are recentered to sum to zero. Dummies use the calendar
quarter, not the row position. After adjustment,
\texttt{gdp\_real\_2015usd} is rebuilt as
\texttt{gdp\_per\_worker $\times$ emp\_persons}, and \texttt{gdp\_sa} is set
to \texttt{SARIMAX}.

Spells shorter than 12 quarters, or missing a calendar quarter, stay
\texttt{NSA}. With the flag off, all 33 IMF-NSA countries remain
\texttt{NSA} (including long spells such as Russia and Sri Lanka).

\section{Exchange rates, including the euro}

$E_{i,2015}$ is the 2015 annual average of domestic currency per USD. For the
United States the value is 1. For the United Kingdom it is the 2015 average
pound per dollar.

Euro-area members as of 1~January~2015---Austria, Belgium, Cyprus, Germany,
Spain, Estonia, Finland, France, Greece, Ireland, Italy, Lithuania,
Luxembourg, Latvia, Malta, the Netherlands, Portugal, Slovakia,
Slovenia---have national \texttt{XDC\_USD} series that \emph{stop at euro
adoption} (Germany: 1998). After adoption the domestic currency is the euro.
Those countries, and euro users without a 2015 national series (Andorra,
Kosovo, Montenegro, San Marino, Vatican), inherit the 2015 Euro Area
aggregate rate, IMF country code \textbf{\texttt{G163}}
($\approx 0.901$ EUR per USD in 2015).

That inherited rate is still a single 2015 converter. Croatia is not in the
set: in 2015 it used the kuna, so its own 2015 \texttt{XDC\_USD} series is the
correct converter. \texttt{G163} is an FX donor only; it is never a panel
country.

\section{Employment}
\label{sec:emp}

The denominator is employed persons, both sexes. ILO \texttt{EMP\_*\_NB}
values are thousands of persons and are multiplied by 1{,}000.

\subsection{Filters}

Applied to every ILO extract:

\begin{itemize}
\item \texttt{sex = SEX\_T} (both sexes). Male/female breakdowns are discarded.
\item Age preference, one definition per country, never mixed:
\begin{enumerate}
\item \texttt{AGE\_YTHADULT\_YGE15} (15+)
\item \texttt{AGE\_AGGREGATE\_YGE15}
\item \texttt{AGE\_AGGREGATE\_TOTAL}
\item \texttt{AGE\_YTHADULT\_Y15-64}
\end{enumerate}
\item One ILO \texttt{source} code per country (the longest series), so LFS
vintages are not spliced.
\end{itemize}

\subsection{Source hierarchy (country-level)}

\begin{enumerate}
\item Official \textbf{quarterly} \texttt{EMP\_TEMP\_SEX\_AGE\_NB\_Q}, if the
country has at least eight quarterly observations after filters.
\item Otherwise official \textbf{annual} \texttt{EMP\_TEMP\_SEX\_AGE\_NB\_A},
linearly interpolated to quarter midpoints.
\item Otherwise ILO modelled annual \texttt{EMP\_2EMP\_SEX\_AGE\_NB\_A},
interpolated the same way.
\end{enumerate}

Sources are not spliced quarter-by-quarter inside a country, except for the
narrow $4\times$ patch in Section~\ref{sec:quality}. In the current vintage,
86 countries use official quarterly employment and 15 use interpolated
official annual. No country falls through to modelled estimates.

\subsection{Annual to quarterly interpolation}

The annual figure for year $y$ is treated as the annual \emph{average} stock,
located at mid-year $t=y+0.5$. Quarter $q$ of year $y$ is located at
\begin{equation}
t = y + (q-0.5)/4
\qquad\text{(Q1: }y+0.125,\;\ldots,\;\text{Q4: }y+0.875\text{)}.
\end{equation}
Between adjacent annual observations the interpolant is the unique straight
line in $t$. Quarters whose midpoint lies outside
$[t_{\mathrm{first}},t_{\mathrm{last}}]$ are dropped (no extrapolation).

This is interpolation of a stock, not Denton benchmarking. The four-quarter
average equals the annual level only if employment is linear in calendar
time. That limitation is accepted in exchange for a transparent, testable
rule.

\section{Data-quality repairs}
\label{sec:quality}

Raw QNEA and ILO extracts contain a small number of unit and survey breaks
that would otherwise dominate the productivity series. The repairs below are
implemented in \texttt{construct.py} and covered by specification tests.

\subsection{Power-of-ten unit breaks in QNEA}

Some QNEA LCU series change unit (1 $\to$ thousands $\to$ millions) without a
matching \texttt{SCALE}. Jamaica in 2018 is a million-fold drop in the raw
extract while ANEA annual real GDP does not drop.

For each country-year compare $4\times\mathrm{mean(quarterly\ LCU)}$ to ANEA
annual real LCU. If they differ by $10^{k}$ with $|k|\ge 3$, that year's
quarterly LCU is multiplied by $10^{-k}$. A SAAR-versus-flow gap of 4
($\log_{10}4\approx 0.6$) is ignored. Years without ANEA inherit the
neighbouring year's factor. The factor is recorded in \texttt{metadata} when
it is not 1.

\subsection{Years that still disagree with ANEA}

After that rescale, country-years whose $4\times\mathrm{mean(quarterly)}$
still differs from ANEA annual real by more than a factor of 5 are
\textbf{dropped}. Residual breaks cannot be repaired with a power of ten
without putting FX back into the volume path.

\subsection{Spells that do not contain the rebase year}

A remaining adjacent-quarter real-GDP move larger than a factor of 20 is
treated as a unit or currency restatement (El Salvador's colon/USD overlap
around dollarization). Only the contiguous spell that contains 2015 is kept,
so the volume index is never a ratio of two different currency units. El
Salvador in the panel therefore starts in 2005Q1 rather than 2000Q1.

\subsection{ILO quarterly employment off the annual path}

Official quarterly employment is occasionally a factor of about four too
small for a calendar year. Israel 2017 is the textbook case: about 0.94
million persons in the quarterly extract against about 3.7 million on the
annual path.

Where an official annual series exists, a quarter whose employment divided
by the interpolated annual path lies outside $[0.4,\,2.5]$ is replaced by
that interpolant and the source string is annotated. Israel 2017 in the
panel is about 3.81 million persons, in line with 2016 and 2018.

\section{Sample and join}

Keep ISO 3166-1 alpha-3 codes (and Kosovo \texttt{KOS} if present). Drop IMF
aggregates. The panel is the \textbf{inner join} of constructed GDP and
employment on $(\mathrm{country},\mathrm{period})$. Duplicate
country--period rows are rejected. Non-positive GDP or employment is
dropped.

The binding constraint is IMF quarterly real GDP, not employment. QNEA
constant-price LCU exists for roughly a hundred economies; the inner join
with 2015 annual levels, 2015 FX, and ILO employment yields 101 countries.

\section{Output schema and metadata}

\begin{table}[h]
\centering
\caption{Columns in \texttt{output/gdp\_per\_worker\_q.csv}.}
\begin{tabular}{lll}
\toprule
Column & Required? & Content \\
\midrule
\texttt{country} & yes & ISO3 \\
\texttt{period} & yes & \texttt{YYYY-Qn} \\
\texttt{gdp\_per\_worker} & yes & annualized 2015 USD per employed person \\
\texttt{metadata} & yes & sources and transformations for the observation \\
\texttt{country\_name} & no & IMF \texttt{CL\_COUNTRY} label \\
\texttt{year}, \texttt{quarter}, \texttt{time} & no & calendar parts; \texttt{time}$=y+(q-1)/4$ \\
\texttt{gdp\_real\_2015usd} & no & annualized real GDP, 2015 USD \\
\texttt{emp\_persons} & no & employed persons \\
\texttt{gdp\_sa} & no & \texttt{SA} or \texttt{NSA} \\
\texttt{emp\_source} & no & ILO series and any interpolation/patch note \\
\bottomrule
\end{tabular}
\end{table}

Each observation's \texttt{metadata} records, in order: the QNEA series and
SA/NSA flag and the $\times 4$ annualization (plus any power-of-ten rescale);
the 2015 real and nominal level sources; the 2015 FX source (including
\texttt{G163} when used); the ILO series, source code, sex, age, and
thousands-to-persons conversion; and that the ratio is 2015 USD market FX,
not PPP.

\section{Replication}

From this folder:
\begin{verbatim}
python download_sources.py
python build_gdp_per_worker.py
python -m unittest tests.test_construct
\end{verbatim}

\texttt{construct.build\_panel(...)} is the construction entry point. The CLI
only reads the cached CSVs into that function. Tests call the shipped
functions (including \texttt{build\_panel}) and check documented identities:
the 2015 USD level, one-for-one pass-through of real LCU growth, euro FX fill
from \texttt{G163}, SA-not-mixed-with-NSA, mid-year interpolation geometry,
no extrapolation, power-of-ten unit alignment, $4\times$ employment patches,
and dropping of IMF staff projections.

\section{Current vintage}

\begin{table}[h]
\centering
\caption{Headline counts, vintage of the \texttt{raw/} extracts dated
10 September 2026.}
\begin{tabular}{lr}
\toprule
Countries & 101 \\
Observations & 7{,}450 \\
Span & 1950Q1--2026Q2 (unbalanced) \\
GDP SA / SARIMAX / NSA (country-level) & 68 / 0 / 33 \\
Employment: official quarterly / interpolated annual & 86 / 15 \\
United States, 2015 mean GDP per worker & $\$122{,}925$ \\
Germany, 2015 mean GDP per worker & $\$85{,}222$ \\
\bottomrule
\end{tabular}
\end{table}

Coverage by country is Table~\ref{tab:coverage}. Longest samples include the
United States (1950Q1--2026Q2, $N=306$), Korea ($N=242$), Canada ($N=202$),
and Australia ($N=192$). Several low-income economies appear with short
spells (Bangladesh 2024 only; Botswana and Mozambique two quarters). Those
short spells are retained: the product is an unbalanced panel, and the
estimator that consumes it already drops countries below a minimum $T$.

\section{Code map}

\begin{tabular}{ll}
\toprule
File & Role \\
\midrule
\texttt{download\_sources.py} & URLs, streaming download, \texttt{raw/} cache \\
\texttt{construct.py} & Every transformation listed above \\
\texttt{build\_gdp\_per\_worker.py} & CLI: download if needed $\to$ build $\to$ \texttt{output/} \\
\texttt{tests/test\_construct.py} & Specification tests against shipped functions \\
\texttt{docs/gdp\_per\_worker\_construction.tex} & This note \\
\bottomrule
\end{tabular}

\section{Citations}

\begin{itemize}
\item IMF, National Economic Accounts, Quarterly (QNEA) and Annual (ANEA).
\item IMF, Exchange Rates (ER), domestic currency per US dollar, period average.
\item ILOSTAT, Employment by sex and age (\texttt{EMP\_TEMP\_SEX\_AGE\_NB});
ILO modelled estimates (\texttt{EMP\_2EMP\_SEX\_AGE\_NB}) as a last-resort
source (unused in this vintage).
\end{itemize}

Cite the vintage of the cached files in \texttt{raw/} (download date $=$ file
timestamps).

\clearpage
\appendix
\section{Country coverage}
% Auto-generated by make_coverage_table.py from output/coverage.csv.
\begin{footnotesize}
\begin{longtable}{l>{\raggedright\arraybackslash}p{2.2in}cccll}
\caption{Country coverage of the constructed panel. Emp.\ is the
country-level employment source: official quarterly ILO (Q official);
official quarterly with isolated $4\times$ patches from interpolated
annual (Q, patched); or linearly interpolated official annual
(A interpolated). SA/NSA is the country-level GDP seasonal-adjustment
choice.}
\label{tab:coverage}\\
\toprule
ISO3 & Country & SA & $N$ & Start & End & Emp. \\
\midrule
\endfirsthead
\toprule
ISO3 & Country & SA & $N$ & Start & End & Emp. \\
\midrule
\endhead
\midrule
\multicolumn{7}{r}{\emph{Continued on next page}}\\
\endfoot
\bottomrule
\endlastfoot
USA & United States & SA & 306 & 1950-Q1 & 2026-Q2 & Q official \\
KOR & Korea, Republic of & SA & 242 & 1966-Q1 & 2026-Q2 & Q official \\
CAN & Canada & SA & 202 & 1976-Q1 & 2026-Q2 & Q official \\
AUS & Australia & SA & 192 & 1978-Q2 & 2026-Q1 & Q official \\
NZL & New Zealand & SA & 156 & 1987-Q2 & 2026-Q1 & Q official \\
NOR & Norway & SA & 139 & 1990-Q1 & 2026-Q1 & Q official \\
SWE & Sweden & SA & 133 & 1993-Q1 & 2026-Q1 & Q official \\
JPN & Japan & SA & 130 & 1994-Q1 & 2026-Q2 & Q official \\
SVK & Slovak Republic & SA & 125 & 1995-Q1 & 2026-Q1 & Q official \\
ESP & Spain & SA & 125 & 1995-Q1 & 2026-Q1 & Q official \\
CZE & Czech Republic & SA & 125 & 1995-Q1 & 2026-Q1 & Q official \\
HUN & Hungary & SA & 122 & 1995-Q1 & 2026-Q1 & Q official \\
PRT & Portugal & SA & 121 & 1995-Q2 & 2026-Q1 & Q official \\
ISR & Israel & SA & 120 & 1995-Q1 & 2024-Q4 & Q official \\
ITA & Italy & SA & 117 & 1996-Q2 & 2026-Q1 & Q official \\
FIN & Finland & SA & 116 & 1995-Q2 & 2026-Q1 & Q official \\
GRC & Greece & SA & 116 & 1995-Q2 & 2026-Q1 & Q official \\
SVN & Slovenia, Republic of & SA & 114 & 1996-Q2 & 2026-Q1 & Q official \\
BEL & Belgium & SA & 113 & 1995-Q2 & 2026-Q1 & Q official \\
FRA & France & SA & 113 & 1983-Q1 & 2026-Q1 & Q official \\
AUT & Austria & SA & 113 & 1995-Q1 & 2026-Q1 & Q official \\
DNK & Denmark & SA & 113 & 1995-Q2 & 2026-Q1 & Q official \\
IRL & Ireland & SA & 112 & 1995-Q2 & 2026-Q1 & Q official \\
ROU & Romania & SA & 111 & 1997-Q2 & 2026-Q1 & Q official \\
MEX & Mexico & SA & 110 & 1995-Q2 & 2026-Q2 & Q official \\
LVA & Latvia, Republic of & SA & 109 & 1996-Q2 & 2026-Q1 & Q official \\
NLD & Netherlands, The & SA & 109 & 1996-Q2 & 2026-Q1 & Q official \\
EST & Estonia, Republic of & SA & 108 & 1997-Q2 & 2026-Q1 & Q official \\
BGR & Bulgaria & SA & 105 & 2000-Q1 & 2026-Q1 & Q official \\
LTU & Lithuania, Republic of & SA & 105 & 1998-Q2 & 2026-Q1 & Q official \\
POL & Poland, Republic of & SA & 105 & 2000-Q1 & 2026-Q1 & Q official \\
HRV & Croatia, Republic of & SA & 105 & 1998-Q1 & 2026-Q1 & Q official \\
MLT & Malta & SA & 103 & 2000-Q2 & 2026-Q1 & Q official \\
HND & Honduras & SA & 102 & 2000-Q1 & 2025-Q2 & A interpolated \\
LUX & Luxembourg & SA & 101 & 1995-Q2 & 2026-Q1 & Q official \\
ISL & Iceland & SA & 100 & 1995-Q2 & 2025-Q4 & Q official \\
PHL & Philippines & SA & 99 & 2000-Q1 & 2024-Q4 & Q official \\
DEU & Germany & SA & 95 & 1991-Q2 & 2026-Q1 & Q official \\
CYP & Cyprus & SA & 93 & 1999-Q2 & 2026-Q1 & Q official \\
MKD & North Macedonia, Republic of & SA & 89 & 2004-Q1 & 2026-Q1 & Q official \\
ZAF & South Africa & SA & 86 & 2000-Q1 & 2025-Q2 & Q official \\
GBR & United Kingdom & SA & 85 & 2005-Q1 & 2026-Q1 & Q official \\
ARG & Argentina & SA & 84 & 2004-Q1 & 2025-Q4 & Q official \\
TUR & T\"urkiye, Republic of & SA & 81 & 2006-Q1 & 2026-Q1 & Q official \\
JAM & Jamaica & SA & 80 & 2000-Q1 & 2023-Q4 & Q official \\
CHE & Switzerland & SA & 79 & 1996-Q2 & 2026-Q1 & Q official \\
SGP & Singapore & SA & 78 & 2000-Q1 & 2025-Q4 & Q official \\
COL & Colombia & SA & 77 & 2007-Q1 & 2026-Q1 & Q official \\
LSO & Lesotho, Kingdom of & SA & 70 & 2007-Q1 & 2024-Q2 & A interpolated \\
KAZ & Kazakhstan, Republic of & SA & 68 & 2003-Q1 & 2019-Q4 & Q official \\
CHL & Chile & SA & 65 & 2010-Q1 & 2026-Q2 & Q official \\
ECU & Ecuador & SA & 64 & 2003-Q4 & 2025-Q3 & Q official \\
CRI & Costa Rica & SA & 63 & 2010-Q3 & 2026-Q1 & Q official \\
SRB & Serbia, Republic of & SA & 60 & 2008-Q2 & 2025-Q4 & Q official \\
MNE & Montenegro & SA & 60 & 2010-Q1 & 2024-Q4 & Q official \\
RUS & Russian Federation & SARIMAX & 60 & 2011-Q1 & 2025-Q4 & Q official \\
LKA & Sri Lanka & SARIMAX & 58 & 2010-Q1 & 2024-Q4 & Q official \\
THA & Thailand & SA & 57 & 2010-Q1 & 2026-Q1 & Q official \\
BRA & Brazil & SA & 57 & 2012-Q1 & 2026-Q1 & Q official \\
ALB & Albania & SA & 52 & 2012-Q1 & 2024-Q4 & Q official \\
TTO & Trinidad and Tobago & SARIMAX & 46 & 2012-Q1 & 2024-Q4 & Q official \\
MDV & Maldives & SARIMAX & 46 & 2003-Q1 & 2014-Q2 & A interpolated \\
DOM & Dominican Republic & SA & 45 & 2015-Q1 & 2026-Q1 & Q official \\
MUS & Mauritius & SARIMAX & 45 & 2013-Q1 & 2024-Q4 & Q official \\
PER & Peru & SARIMAX & 44 & 2011-Q1 & 2021-Q4 & Q official \\
BRN & Brunei Darussalam & SARIMAX & 44 & 2014-Q3 & 2025-Q2 & A interpolated \\
MNG & Mongolia & SARIMAX & 44 & 2015-Q1 & 2025-Q4 & Q official \\
UKR & Ukraine & SA & 42 & 2010-Q1 & 2020-Q2 & Q official \\
CPV & Cabo Verde & SARIMAX & 40 & 2009-Q3 & 2019-Q2 & A interpolated \\
GMB & Gambia, The & SARIMAX & 40 & 2014-Q1 & 2023-Q4 & A interpolated \\
KOS & Kosovo, Republic of & SARIMAX & 40 & 2012-Q1 & 2021-Q4 & Q official \\
BHS & Bahamas, The & SARIMAX & 38 & 2015-Q1 & 2024-Q2 & A interpolated \\
PRY & Paraguay & SARIMAX & 36 & 2017-Q1 & 2025-Q4 & Q official \\
IND & India & SA & 34 & 2017-Q3 & 2025-Q4 & Q official \\
IDN & Indonesia & SA & 32 & 2008-Q1 & 2023-Q3 & Q official \\
NIC & Nicaragua & SARIMAX & 30 & 2006-Q1 & 2013-Q2 & A interpolated \\
RWA & Rwanda & SARIMAX & 28 & 2017-Q1 & 2025-Q4 & Q official \\
LCA & St. Lucia & SARIMAX & 28 & 2017-Q1 & 2024-Q3 & Q official \\
ARM & Armenia, Republic of & SARIMAX & 24 & 2008-Q1 & 2013-Q4 & Q official \\
SYC & Seychelles & SARIMAX & 23 & 2016-Q1 & 2024-Q2 & Q, patched \\
NGA & Nigeria & SARIMAX & 23 & 2014-Q1 & 2024-Q2 & Q official \\
SEN & Senegal & SA & 23 & 2016-Q2 & 2025-Q1 & Q official \\
BIH & Bosnia and Herzegovina & SA & 23 & 2020-Q1 & 2025-Q3 & Q official \\
GEO & Georgia & SARIMAX & 23 & 2020-Q1 & 2025-Q3 & Q official \\
NAM & Namibia & SARIMAX & 22 & 2013-Q1 & 2018-Q2 & A interpolated \\
EGY & Egypt, Arab Republic of & SARIMAX & 21 & 2019-Q4 & 2024-Q4 & Q official \\
ZMB & Zambia & SARIMAX & 20 & 2017-Q1 & 2024-Q4 & Q official \\
GTM & Guatemala & SARIMAX & 20 & 2013-Q2 & 2025-Q3 & Q official \\
CHN & China, People's Republic of & SARIMAX & 18 & 2011-Q1 & 2015-Q2 & A interpolated \\
QAT & Qatar & SARIMAX & 16 & 2013-Q1 & 2016-Q4 & Q official \\
PAK & Pakistan & SARIMAX & 16 & 2017-Q3 & 2025-Q4 & Q official \\
MAR & Morocco & SA & 15 & 2014-Q1 & 2017-Q3 & Q official \\
UGA & Uganda & SA & 12 & 2016-Q3 & 2019-Q2 & A interpolated \\
GHA & Ghana & NSA & 11 & 2022-Q1 & 2024-Q3 & Q official \\
SLV & El Salvador & NSA & 10 & 2005-Q1 & 2007-Q2 & A interpolated \\
SAU & Saudi Arabia & SA & 8 & 2012-Q1 & 2016-Q1 & Q official \\
KEN & Kenya & SA & 8 & 2019-Q1 & 2021-Q4 & Q official \\
BLR & Belarus, Republic of & NSA & 6 & 2014-Q1 & 2015-Q2 & A interpolated \\
BGD & Bangladesh & NSA & 4 & 2024-Q1 & 2024-Q4 & Q official \\
BWA & Botswana & NSA & 2 & 2006-Q1 & 2006-Q2 & A interpolated \\
MOZ & Mozambique, Republic of & NSA & 2 & 2022-Q1 & 2022-Q2 & A interpolated \\
\end{longtable}
\end{footnotesize}


\end{document}
